
\documentclass{article}
\usepackage{colortbl}
\usepackage{makecell}
\usepackage{multirow}
\usepackage{supertabular}

\begin{document}

\newcounter{utterance}

\centering \large Interaction Transcript for game `wordle', experiment `medium\_frequency\_words\_no\_clue\_no\_critic', episode 11 with reflexion{-}Qwen3.5{-}27B.
\vspace{24pt}

{ \footnotesize  \setcounter{utterance}{1}
\setlength{\tabcolsep}{0pt}
\begin{supertabular}{c@{$\;$}|p{.15\linewidth}@{}p{.15\linewidth}p{.15\linewidth}p{.15\linewidth}p{.15\linewidth}p{.15\linewidth}}
    \# & \multicolumn{2}{c}{Player} && \multicolumn{2}{c}{Game Master} \\
    \hline

    \theutterance \stepcounter{utterance}  
    & & & \multicolumn{4}{p{0.6\linewidth}}{
        \cellcolor[rgb]{0.9,0.9,0.9}{
            \makecell[{{p{\linewidth}}}]{
                \texttt{\tiny{[P1$\langle$GM]}}
                \texttt{You are a language wizard who likes to guess words by using the given rules.} \\
\\ 
\texttt{Welcome to Wordle! You have six attempts to guess the target word, a valid English word of five lowercase letters (a{-}z). Please use the tags "explanation:" and "guess:" to provide a concise explanation for each guess.} \\
\\ 
\texttt{For instance, if your guess is "apple", your response should be} \\
\texttt{explanation: this is a common five{-}letter English word, and I am starting my guess with this word.} \\
\texttt{guess: apple} \\
\\ 
\texttt{After each guess, your answer will be validated, and you will receive feedback indicating which letters are correct (green), which letters are correct but in the wrong position (yellow), and which letters are incorrect (red). This feedback can be useful in determining which letters to include or exclude in your next guess.} \\
\\ 
\texttt{For example, the feedback for "apple" might be:} \\
\texttt{guess\_feedback: a<yellow> p<yellow> p<green> l<yellow> e<red>} \\
\\ 
\texttt{The explanation should contain details about how the guess\_feedback is used to arrive at a new guess.} \\
\\ 
\texttt{Do not generate any other text.} \\
\\ 
\texttt{Let's begin with your first guess.} \\
            }
        }
    }
     \\ \\

    \theutterance \stepcounter{utterance}  
    & \multicolumn{4}{p{0.6\linewidth}}{
        \cellcolor[rgb]{0.561,0.737,0.561}{
            \makecell[{{p{\linewidth}}}]{
                \texttt{\tiny{[P1$\rangle$GM]}}
                \texttt{explanation: I am starting with a common five{-}letter word containing frequent vowels and consonants to maximize the information gained from the first feedback.} \\
\texttt{guess: slate} \\
            }
        }
    }
    & & \\ \\

    \theutterance \stepcounter{utterance}  
    & & & \multicolumn{2}{p{0.3\linewidth}}{
        \cellcolor[rgb]{0.9,0.9,0.9}{
            \makecell[{{p{\linewidth}}}]{
                \texttt{\tiny{[GM$|$GM]}}
                \texttt{attempts = 1} \\
\texttt{target\_word = faint} \\
\texttt{guess = slate} \\
\texttt{guess\_feedback = s<red> l<red> a<yellow> t<yellow> e<red>} \\
            }
        }
    }
    & & \\ \\

    \theutterance \stepcounter{utterance}  
    & & & \multicolumn{4}{p{0.6\linewidth}}{
        \cellcolor[rgb]{0.9,0.9,0.9}{
            \makecell[{{p{\linewidth}}}]{
                \texttt{\tiny{[P1$\langle$GM]}}
                \texttt{guess\_feedback: s<red> l<red> a<yellow> t<yellow> e<red>} \\
\\ 
\texttt{Provide your response only in this format.} \\
\texttt{explanation: details} \\
\texttt{guess: word} \\
            }
        }
    }
     \\ \\

    \theutterance \stepcounter{utterance}  
    & \multicolumn{4}{p{0.6\linewidth}}{
        \cellcolor[rgb]{0.561,0.737,0.561}{
            \makecell[{{p{\linewidth}}}]{
                \texttt{\tiny{[P1$\rangle$GM]}}
                \texttt{explanation: The feedback indicates 'a' and 't' are in the word but in the wrong positions (yellow), while 's', 'l', and 'e' are not in the word (red). I need a word containing 'a' and 't' but avoiding 's', 'l', and 'e', placing 'a' and 't' in new positions. "Crate" fits these criteria, moving 'a' to the second position and 't' to the fourth, while introducing new common letters 'c' and 'r'.} \\
\texttt{guess: crate} \\
            }
        }
    }
    & & \\ \\

    \theutterance \stepcounter{utterance}  
    & & & \multicolumn{2}{p{0.3\linewidth}}{
        \cellcolor[rgb]{0.9,0.9,0.9}{
            \makecell[{{p{\linewidth}}}]{
                \texttt{\tiny{[GM$|$GM]}}
                \texttt{attempts = 2} \\
\texttt{target\_word = faint} \\
\texttt{guess = crate} \\
\texttt{guess\_feedback = c<red> r<red> a<yellow> t<yellow> e<red>} \\
            }
        }
    }
    & & \\ \\

    \theutterance \stepcounter{utterance}  
    & & & \multicolumn{4}{p{0.6\linewidth}}{
        \cellcolor[rgb]{0.9,0.9,0.9}{
            \makecell[{{p{\linewidth}}}]{
                \texttt{\tiny{[P1$\langle$GM]}}
                \texttt{guess\_feedback: c<red> r<red> a<yellow> t<yellow> e<red>} \\
\\ 
\texttt{Provide your response only in this format.} \\
\texttt{explanation: details} \\
\texttt{guess: word} \\
            }
        }
    }
     \\ \\

    \theutterance \stepcounter{utterance}  
    & \multicolumn{4}{p{0.6\linewidth}}{
        \cellcolor[rgb]{0.561,0.737,0.561}{
            \makecell[{{p{\linewidth}}}]{
                \texttt{\tiny{[P1$\rangle$GM]}}
                \texttt{explanation: The feedback confirms 'a' and 't' are in the word but in the wrong positions (yellow), while 's', 'l', 'e', 'c', and 'r' are not in the word (red). I need a word with 'a' and 't' avoiding the excluded letters. Since 'a' was yellow in position 2 (slate) and 3 (crate), and 't' was yellow in position 4 (slate) and 4 (crate), I will try placing 'a' in position 1 and 't' in position 5, using common remaining consonants like 'p', 'i', 'n'. "Pinta" fits these constraints.} \\
\texttt{guess: pinta} \\
            }
        }
    }
    & & \\ \\

    \theutterance \stepcounter{utterance}  
    & & & \multicolumn{2}{p{0.3\linewidth}}{
        \cellcolor[rgb]{0.9,0.9,0.9}{
            \makecell[{{p{\linewidth}}}]{
                \texttt{\tiny{[GM$|$GM]}}
                \texttt{attempts = 3} \\
\texttt{target\_word = faint} \\
\texttt{guess = pinta} \\
\texttt{guess\_feedback = p<red> i<yellow> n<yellow> t<yellow> a<yellow>} \\
            }
        }
    }
    & & \\ \\

    \theutterance \stepcounter{utterance}  
    & & & \multicolumn{4}{p{0.6\linewidth}}{
        \cellcolor[rgb]{0.9,0.9,0.9}{
            \makecell[{{p{\linewidth}}}]{
                \texttt{\tiny{[P1$\langle$GM]}}
                \texttt{guess\_feedback: p<red> i<yellow> n<yellow> t<yellow> a<yellow>} \\
\\ 
\texttt{Provide your response only in this format.} \\
\texttt{explanation: details} \\
\texttt{guess: word} \\
            }
        }
    }
     \\ \\

    \theutterance \stepcounter{utterance}  
    & \multicolumn{4}{p{0.6\linewidth}}{
        \cellcolor[rgb]{0.561,0.737,0.561}{
            \makecell[{{p{\linewidth}}}]{
                \texttt{\tiny{[P1$\rangle$GM]}}
                \texttt{explanation: The feedback shows 'p' is incorrect (red), while 'i', 'n', 't', and 'a' are all in the word but in the wrong positions (yellow). From previous guesses, 's', 'l', 'e', 'c', 'r' are also excluded. The word must contain \{i, n, t, a\} plus one unknown letter. 'a' was yellow in pos 2, 3, 1; 't' was yellow in pos 4, 4, 4; 'i' was yellow in pos 2; 'n' was yellow in pos 3. I need to rearrange these. If I try "taint", 't' is pos 1 (new), 'a' is pos 2 (was yellow in 2, 3, 1 {-} wait, 'a' was yellow in pos 2 in 'slate', pos 3 in 'crate', pos 1 in 'pinta'. So 'a' cannot be in 1, 2, or 3. 'a' must be in 4 or 5. 't' was yellow in 4, 4, 4. So 't' cannot be in 4. 't' must be in 1, 2, 3, or 5. 'i' was yellow in 2. 'n' was yellow in 3. Let's try "tania" (t{-}1, a{-}2, n{-}3, i{-}4, a{-}5) {-} wait, 'a' cannot be 2. Let's try "tina" + something. How about "taint"? t(1), a(2), i(3), n(4), t(5). 'a' at 2 is bad. 'n' at 4 is new. 'i' at 3 is new. 't' at 1 is new. 't' at 5 is new. But 'a' at 2 was yellow in 'slate'. So 'a' != 2. Let's try "tania" again? No. Let's try "titan"? t(1), i(2), t(3), a(4), n(5). 'i' at 2 was yellow in 'pinta'. So 'i' != 2. 'a' at 4 is new. 'n' at 5 is new. 't' at 1, 3. 't' was yellow in 4. So 1, 3 are okay. 'i' was yellow in 2. So 2 is bad. 'n' was yellow in 3. So 3 is bad. 'a' was yellow in 1, 2, 3. So 1, 2,} \\
            }
        }
    }
    & & \\ \\

    \theutterance \stepcounter{utterance}  
    & & & \multicolumn{2}{p{0.3\linewidth}}{
        \cellcolor[rgb]{0.9,0.9,0.9}{
            \makecell[{{p{\linewidth}}}]{
                \texttt{\tiny{[GM$|$GM]}}
                \texttt{Error: The response should contain the 'guess:' keyword exactly once.} \\
            }
        }
    }
    & & \\ \\

    \theutterance \stepcounter{utterance}  
    & & & \multicolumn{2}{p{0.3\linewidth}}{
        \cellcolor[rgb]{0.9,0.9,0.9}{
            \makecell[{{p{\linewidth}}}]{
                \texttt{\tiny{[GM$|$GM]}}
                \texttt{game\_result = ABORT} \\
            }
        }
    }
    & & \\ \\

    \theutterance \stepcounter{utterance}  
    & & \multicolumn{4}{p{0.6\linewidth}}{
        \cellcolor[rgb]{0.95,0.95,0.95}{
            \makecell[{{p{\linewidth}}}]{
                \texttt{\tiny{[P1$\leftrightarrow$P]}}
                \texttt{Post{-}episode Reflexion (not sent to GM):} \\
\texttt{New plan: The previous attempt failed because I included my internal reasoning process (e.g., "wait, 'a' was yellow...", "Let's try...") directly in the "explanation" field, causing the response to become disorganized and likely violating the "concise explanation" rule or the strict format. Additionally, I did not provide a final "guess:" line before the text cut off. Next time, I will strictly separate my internal logic from the output. I will perform the deduction mentally to identify the correct word (likely "taint" or "tania" based on the constraints: 'a' not in 1,2,3; 't' not in 4; 'i' not in 2; 'n' not in 3), and then output *only* a clean, concise explanation of the final choice followed immediately by the "guess:" tag and the five{-}letter word. I will ensure the explanation summarizes the logic without showing the step{-}by{-}step trial{-}and{-}error or self{-}correction.} \\
            }
        }
    }
    & \\ \\

\end{supertabular}
}

\end{document}
